\documentclass[../../main.tex]{subfiles}
\begin{document}
% SOURCE p=241; printed=226; object=ch07-radon-products-title; region=top
\section{Radon 测度的乘积}\label{section:ch07-radon-products}\label{section:ch07-products-topological-spaces}
% SOURCE p=241; printed=226; object=ch07-radon-products-introduction; region=upper
本节研究乘积空间上的 Radon 测度。$X,Y$ 表示 LCH 空间，$\pi_X,\pi_Y$ 分别表示 $X\times Y$ 到 $X,Y$ 的投影。
% SOURCE p=241; printed=226; object=ch07-product-borel-radon; number=7.20; region=upper
\begin{theorem}\label{theorem:ch07-product-borel-radon}
\begin{enumerate}
\item\label{item:ch07-product-borel-inclusion} $\mathscr B_X\otimes\mathscr B_Y\subseteq\mathscr B_{X\times Y}$。
\item\label{item:ch07-product-borel-equality} 若 $X,Y$ 都第二可数，则 $\mathscr B_X\otimes\mathscr B_Y=\mathscr B_{X\times Y}$。
\item\label{item:ch07-second-countable-product-radon} 若 $X,Y$ 都第二可数，且 $\mu,\nu$ 分别为 $X,Y$ 上的 Radon 测度，则 $\mu\times\nu$ 是 $X\times Y$ 上的 Radon 测度。
\end{enumerate}
\end{theorem}
% SOURCE p=241; printed=226; object=ch07-product-borel-radon-proof; region=middle
\begin{proof}
前两项直接推广了\cref{proposition:ch01-product-borel}，证明本质相同。主要工具是\cref{proposition:ch01-product-generating-families}：$\mathscr B_X\otimes\mathscr B_Y$ 由 $U\times V$ 型集合生成，其中 $U$ 在 $X$ 中开、$V$ 在 $Y$ 中开。这些矩形在乘积空间中开，故得到第一项。

若两个拓扑有可数基 $\mathcal E,\mathcal F$，则 $X,Y,X\times Y$ 中的每个开集分别是 $\mathcal E,\mathcal F,\{U\times V:U\in\mathcal E,V\in\mathcal F\}$ 中集合的可数并。因此相应 Borel 代数由这些族生成，得到第二项。至于第三项，由第二项，$\mu\times\nu$ 是 Borel 测度；根据\cref{theorem:ch07-automatic-regularity}，只需证明每个紧集 $K\subseteq X\times Y$ 测度有限。两投影 $\pi_X(K),\pi_Y(K)$ 都紧，且 $K\subseteq\pi_X(K)\times\pi_Y(K)$，故
\begin{align*}
(\mu\times\nu)(K)\leqslant\mu(\pi_X(K))\nu(\pi_Y(K))<\infty.
\end{align*}
原书在这里把 $\leqslant$ 印作了等号；所需有限性由不等式即可得到。
\end{proof}
% SUPPLEMENT: id=v14-ch07-product-borel-radon-details; parent=ch07-product-borel-radon; kind=proof-detail+correction+motivation+hypotheses; origin=independent-derivation
\begin{remark}
\textbf{动机、条件与末式订正（重构新增）。}乘积 $\sigma$ 代数由矩形经\emph{可数}操作生成，乘积拓扑则允许任意并，因而二者可能不同。可数基使每个开集只需可数个基矩形，正好消除差别。第一、二项本身不需 LCH；第三项用 LCH 与第二可数保证自动正则性。有限个 LCH 空间的乘积仍为 LCH：两个紧邻域的乘积是紧邻域，两个 Hausdorff 空间的乘积仍 Hausdorff。

原末式为 $(\mu\times\nu)(K)=\mu(\pi_X(K))\nu(\pi_Y(K))$，一般不成立。取 $X=Y=\mathbb R$、两测度为 Lebesgue 测度，$K=\{(t,t):0\leqslant t\leqslant1\}$。两投影质量都是 $1$，但 Fubini 定理给 $K$ 的二维测度为零，因为每个纵截面至多一个点。正确的不等式来自单调性及矩形质量公式。两个投影质量有限，故这里不存在 $0\cdot\infty$ 的歧义；空集时两边为零。
\end{remark}
% END SUPPLEMENT: id=v14-ch07-product-borel-radon-details
% SOURCE p=241; printed=226; object=ch07-product-borel-obstruction; region=middle-lower
当 $X$ 或 $Y$ 不是第二可数时，可能有 $\mathscr B_X\otimes\mathscr B_Y\ne\mathscr B_{X\times Y}$，见\cref{exercise:ch07-ordinal-product-borel,exercise:ch07-large-discrete-product-borel}。这时 Radon 测度的普通乘积当然不是整个乘积空间上的 Radon 测度。不过，可以从它自然地构造一个 Radon 测度。为说明这一点，先需要两个关于连续函数的事实。
% SOURCE p=241; printed=226; object=ch07-tensor-function-definition; region=middle-lower
\begin{definition}[函数的张量积]\label{definition:ch07-tensor-functions}
若 $g,h$ 分别为 $X,Y$ 上的函数，定义 $X\times Y$ 上的函数 $g\otimes h$ 为
\begin{align*}
(g\otimes h)(x,y)=g(x)h(y).
\end{align*}
\end{definition}
% SUPPLEMENT: id=v14-ch07-tensor-functions-details; parent=ch07-tensor-function-definition; kind=motivation+hypotheses; origin=independent-derivation
\begin{remark}
\textbf{动机与约定（重构新增）。}这里的张量积先指一个具体的乘积函数；它把两个变量分开，使积分可以逐个变量计算。若两函数连续，则乘积连续；若都有紧支集，则乘积支集包含于两支集的紧乘积，故属于 $C_c(X\times Y)$。前段所说“不是 Radon 测度”首先是定义域问题：普通乘积只在乘积 $\sigma$ 代数上定义，而本章的 Radon 测度须定义在全部 Borel 集上；并不否认它可能有合适的延拓。
\end{remark}
% END SUPPLEMENT: id=v14-ch07-tensor-functions-details
% SOURCE p=241; printed=226; object=ch07-tensor-density; number=7.21; region=lower
\begin{proposition}\label{proposition:ch07-tensor-density}
令 $\mathcal P$ 为所有 $g\otimes h$（$g\in C_c(X),h\in C_c(Y)$）张成的向量空间。则 $\mathcal P$ 关于一致范数在 $C_c(X\times Y)$ 中稠密。更确切地，给定 $f\in C_c(X\times Y)$、$\varepsilon>0$，以及分别包含 $\pi_X(\operatorname{supp}f)$、$\pi_Y(\operatorname{supp}f)$ 的相对紧开集 $U\subseteq X,V\subseteq Y$，存在 $F\in\mathcal P$，使 $\|F-f\|_u<\varepsilon$ 且 $\operatorname{supp}F\subseteq U\times V$。
\end{proposition}
% SOURCE p=241; printed=226; object=ch07-tensor-density-proof; region=bottom
\begin{proof}
$\overline U\times\overline V$ 是紧 Hausdorff 空间。由 Stone--Weierstrass 定理容易得到，$\{g\otimes h:g\in C(\overline U),h\in C(\overline V)\}$ 的线性张成空间在 $C(\overline U\times\overline V)$ 中稠密。因此其中有 $G$，使 $\sup_{\overline U\times\overline V}|G-f|<\varepsilon$。又由 Urysohn 引理，存在 $\phi\in C_c(U,[0,1])$、$\psi\in C_c(V,[0,1])$，分别在两个支集投影上等于 $1$。在 $\overline U\times\overline V$ 上定义 $F=(\phi\otimes\psi)G$，在其他地方定义 $F=0$，则 $F\in\mathcal P$，$\operatorname{supp}F\subseteq U\times V$，且 $\|F-f\|_u<\varepsilon$。
\end{proof}
% SUPPLEMENT: id=v14-ch07-tensor-density-details; parent=ch07-tensor-density; kind=proof-detail+motivation+hypotheses; origin=independent-derivation
\begin{remark}
\textbf{逼近与零延拓的核验（重构新增）。}紧乘积上的上述线性张成空间是代数，含常数并对共轭封闭。两点若在 $X$ 坐标不同，可用紧 Hausdorff 空间 $\overline U$ 上分离该两点的连续函数及 $Y$ 上常数 $1$ 分离；$Y$ 坐标不同也一样。因此复 Stone--Weierstrass 定理适用。

写 $G=\sum_{j=1}^mg_j\otimes h_j$。每个 $\phi g_j$ 在 $U$ 上连续，支集包含于紧集 $\operatorname{supp}\phi\subseteq U$；在 $X\setminus U$ 上延成零仍连续，因为支集外有开邻域上恒为零。它于是属于 $C_c(X)$；$\psi h_j$ 同理。因此零延拓后的 $F=\sum_j(\phi g_j)\otimes(\psi h_j)$ 确属 $\mathcal P$。在整个乘积上有 $(\phi\otimes\psi)f=f$，所以矩形内 $|F-f|=|\phi\otimes\psi|\,|G-f|<\varepsilon$，矩形外二者均为零。其支集包含于 $\operatorname{supp}\phi\times\operatorname{supp}\psi\subseteq U\times V$。若 $f=0$，直接取 $F=0$ 即可。

动机是以有限个分离变量函数逼近一般连续函数，同时保持一个固定的有限测度容器。相对紧性提供紧闭包以应用逼近定理，LCH 提供紧支撑截断；只有局部一致逼近而没有统一支集控制，后面的积分误差未必可控。
\end{remark}
% END SUPPLEMENT: id=v14-ch07-tensor-density-details
% SOURCE p=242; printed=227; object=ch07-cc-product-integrals; number=7.22; region=top
\begin{proposition}\label{proposition:ch07-cc-product-integrals}
每个 $f\in C_c(X\times Y)$ 都是 $\mathscr B_X\otimes\mathscr B_Y$ 可测的。此外，若 $\mu,\nu$ 分别为 $X,Y$ 上的 Radon 测度，则 $C_c(X\times Y)\subseteq L^1(\mu\times\nu)$，且
\begin{align*}
\int f\,\mathrm d(\mu\times\nu)=\iint f\,\mathrm d\mu\,\mathrm d\nu
=\iint f\,\mathrm d\nu\,\mathrm d\mu
\qquad(f\in C_c(X\times Y)).
\end{align*}
\end{proposition}
% SOURCE p=242; printed=227; object=ch07-cc-product-integrals-proof; region=upper
\begin{proof}
若 $g\in C_c(X),h\in C_c(Y)$，则 $g\otimes h=(g\circ\pi_X)(h\circ\pi_Y)$。由乘积 $\sigma$ 代数的定义，两个投影分别可测到 $\mathscr B_X,\mathscr B_Y$；而 $g,h$ 连续，所以这两个复合函数可测。可测函数的积、和、逐点极限仍可测，结合\cref{proposition:ch07-tensor-density}便得第一项。此外，每个 $f$ 有界，且支撑于一个有限乘积测度的集合，故可积。即使 $\mu,\nu$ 不为 $\sigma$ 有限，Fubini 定理对这种 $f$ 仍成立，因为可以把它们换成在 $\pi_X(\operatorname{supp}f)$ 与 $\pi_Y(\operatorname{supp}f)$ 上的有限测度限制。
\end{proof}
% SUPPLEMENT: id=v14-ch07-cc-product-integrals-details; parent=ch07-cc-product-integrals; kind=proof-detail+motivation+hypotheses; origin=independent-derivation
\begin{remark}
\textbf{局部化的完整说明（重构新增）。}取误差 $1/n$ 的张量和近似，得到一个逐点收敛序列，因而确可使用可测函数序列极限定理。令 $K_X=\pi_X(\operatorname{supp}f)$、$K_Y=\pi_Y(\operatorname{supp}f)$。它们是紧 Borel 集，矩形 $K_X\times K_Y$ 属于乘积 $\sigma$ 代数，且测度等于 $\mu(K_X)\nu(K_Y)<\infty$。于是
\begin{align*}
\int|f|\,\mathrm d(\mu\times\nu)\leqslant\|f\|_u\mu(K_X)\nu(K_Y)<\infty.
\end{align*}
这里使用可测的紧矩形作为容器，不需要先断言任意紧子集本身都在乘积 $\sigma$ 代数内。

普通乘积测度在这个矩形上的限制，等于有限限制测度 $\mu|_{K_X}$ 与 $\nu|_{K_Y}$ 的乘积：任意矩形覆盖与 $K_X\times K_Y$ 相交，仍为矩形覆盖且成本不增，反向局部矩形覆盖也可看作全空间的覆盖，所以定义外测度的下确界相同。对两个有限空间应用 Fubini 定理，得到两种积分次序相等。$K_X$ 或 $K_Y$ 外的相应截面恒为零，所以局部积分与全空间的累次积分相同。这说明紧支撑用有限局部化替代了全局 $\sigma$ 有限性，而不是直接删去了普通 Fubini 定理的假设。
\end{remark}
% END SUPPLEMENT: id=v14-ch07-cc-product-integrals-details
\subfile{section-04-radon-product.tex}
\end{document}
